\documentclass[11pt,a4paper]{article}

% ── Page geometry ─────────────────────────────────────────────────────────────
\usepackage[a4paper, top=2.5cm, bottom=2.5cm, left=2.5cm, right=2.5cm]{geometry}

% ── Fonts & encoding ──────────────────────────────────────────────────────────
\usepackage[T1]{fontenc}
\usepackage[utf8]{inputenc}
\usepackage{lmodern}
\usepackage{microtype}

% ── Hyperlinks ─────────────────────────────────────────────────────────────────
\usepackage[colorlinks=true, linkcolor=black, citecolor=black, urlcolor=blue]{hyperref}

% ── Math ──────────────────────────────────────────────────────────────────────
\usepackage{amsmath, amssymb}

% ── Tables ────────────────────────────────────────────────────────────────────
\usepackage{booktabs}
\usepackage{array}

% ── Code listings ─────────────────────────────────────────────────────────────
\usepackage{listings}
\usepackage{xcolor}

\definecolor{codebg}{rgb}{0.96,0.96,0.96}
\definecolor{codeframe}{rgb}{0.80,0.80,0.80}

\lstset{
  backgroundcolor=\color{codebg},
  frame=single,
  rulecolor=\color{codeframe},
  basicstyle=\small\ttfamily,
  breaklines=true,
  captionpos=b,
  xleftmargin=1em,
  xrightmargin=1em,
  aboveskip=0.7em,
  belowskip=0.7em,
  showstringspaces=false
}

% ── Section spacing ───────────────────────────────────────────────────────────
\usepackage{titlesec}
\titlespacing*{\section}{0pt}{1.3ex plus 0.4ex}{0.7ex}
\titlespacing*{\subsection}{0pt}{1.0ex plus 0.3ex}{0.5ex}

% ── Paragraph spacing ─────────────────────────────────────────────────────────
\setlength{\parskip}{0.4em}
\setlength{\parindent}{0pt}

% ── Abstract width ────────────────────────────────────────────────────────────
\usepackage{abstract}
\renewcommand{\abstractnamefont}{\normalfont\bfseries}

% ─────────────────────────────────────────────────────────────────────────────
\begin{document}

% ── Title block ───────────────────────────────────────────────────────────────
\title{\textbf{QUANTA: Engineering a Production-Ready\\
Post-Quantum Blockchain with Falcon-512 Lattice Signatures}}

\author{
  Kishore K\\
  \small Independent Researcher\\
  \small \href{mailto:kishorekkumar34@gmail.com}{kishorekkumar34@gmail.com}\\
  \small \url{https://www.quantachain.org} \quad
  \url{https://github.com/quantachain/quanta}
}

\date{February 2026}
\maketitle
\thispagestyle{empty}

% ── Abstract ──────────────────────────────────────────────────────────────────
\begin{abstract}
We present QUANTA, an open-source blockchain built exclusively on
NIST-standardized post-quantum cryptography: Falcon-512 (FIPS~206) for
transaction signatures and Kyber-1024 (ML-KEM, FIPS~203) for node-to-node
key encapsulation. No classical primitive appears in the critical path.
We describe five engineering contributions: (1)~a domain-separated
Falcon-512 verification pipeline eliminating malleability and cross-protocol
replay; (2)~a parallel signature-verification engine achieving a 6$\times$
speedup on commodity multi-core hardware; (3)~an LRU-backed signature cache
with $\approx$80\% hit rates; (4)~a binary+Zstandard storage layer reducing
on-disk block size by 75\% vs.\ plain JSON; and (5)~a single-byte
crypto-agility tag enabling fork-free algorithm migration. We also give
informal security arguments for five attack classes. QUANTA is MIT-licensed;
a pre-mainnet testnet is publicly accessible via Docker.
Extended version: \href{https://doi.org/10.5281/zenodo.18753528}{doi:10.5281/zenodo.18753528}.
\end{abstract}

\hrule
\vspace{0.8em}

% ─────────────────────────────────────────────────────────────────────────────
\section{Introduction}

Deployed blockchains rely almost entirely on the hardness of the elliptic
curve discrete logarithm problem (ECDLP). Both ECDSA and EdDSA are broken
by Shor's algorithm~\cite{shor1994} on a cryptographically relevant quantum
computer (CRQC). Conservative projections place CRQCs capable of attacking
256-bit elliptic curve keys within 10--15 years~\cite{mosca2018}. Because
blockchain public keys are permanently on-chain, a harvest-now/decrypt-later
attacker can record public keys today and forge signatures retroactively.
The blockchain setting is acute: once a CRQC arrives, every address that has
ever broadcast a transaction has an exposed public key and is permanently
vulnerable.

Existing post-quantum (PQ) blockchain efforts fall into two categories:
hybrid retrofits, where classical keys coexist with PQ keys during a
transition period (e.g., Algorand's announced migration research), and
purpose-built PQ chains, the most notable being the Quantum Resistant
Ledger (QRL)~\cite{hulsing2018}, which uses XMSS---a stateful hash-based
scheme that requires careful per-key signing-state management.

QUANTA belongs to the second category but distinguishes itself in three
ways. First, it uses Falcon-512~\cite{falcon2020}, a NIST-standardized
(FIPS~206) stateless lattice scheme, avoiding the key-state problem of XMSS
while providing compact keys and signatures. Second, it solves the
determinism problem that arises when variable-time Falcon operations interact
with multi-architecture consensus: all arithmetic is pure integer math and
the compiler toolchain is pinned (\S\ref{sec:determinism}). Third, it pairs
Falcon-512 with Kyber-1024 (ML-KEM, FIPS~203) for P2P confidentiality,
making both authentication and transport layers quantum-resistant from
genesis---a combination not present in any prior PQ blockchain system.

\begin{table}[h]
\centering
\small
\begin{tabular}{@{}lcccc@{}}
\toprule
\textbf{Property} & \textbf{QRL} & \textbf{QAN} & \textbf{Algorand PQ} & \textbf{QUANTA} \\
\midrule
Signature type    & XMSS (stateful) & Dilithium & ECDSA+Falcon hybrid & Falcon-512 (stateless) \\
KEM layer         & None & None & None & Kyber-1024 \\
Consensus determinism & -- & -- & -- & Explicit (\S\ref{sec:determinism}) \\
Classical fallback & No  & Yes  & Yes  & No \\
Compressed storage & No  & No   & No   & Yes (75\%) \\
Crypto-agility    & No  & No   & No   & Yes (tag-dispatch) \\
\bottomrule
\end{tabular}
\caption{Comparison with prior post-quantum blockchain systems.}
\label{tab:comparison}
\end{table}

Hybrid approaches retain classical keys in the critical path; by the
weakest-link argument, a sufficiently powerful CRQC can still compromise
hybrid-signed transactions by attacking the classical component. QUANTA
removes all classical primitives at genesis.

\textbf{Contributions.}
\begin{itemize}
  \item A \textbf{protocol hardening framework} for Falcon-512 in consensus, with informal security arguments for five attack classes (\S\ref{sec:hardening}).
  \item A \textbf{parallel and cached block-validation engine} projected at $6{\times}$ speedup (\S\ref{sec:parallel}).
  \item A \textbf{determinism enforcement model} covering arithmetic, toolchain, and compilation (\S\ref{sec:determinism}).
  \item A \textbf{compressed storage architecture} reducing disk cost by 75\% (\S\ref{sec:storage}).
  \item A \textbf{crypto-agility mechanism} enabling fork-free algorithm migration (\S\ref{sec:agility}).
\end{itemize}

% ─────────────────────────────────────────────────────────────────────────────
\section{Background}

\subsection{Falcon-512 and Kyber-1024}

Falcon~\cite{falcon2020} is a lattice-based signature scheme based on NTRU
lattices and Fast Fourier sampling, standardized by NIST as FIPS~206. At
security level~1 (Falcon-512): public key 897\,B, secret key
$\approx$1281\,B, signature 33--666\,B (variable, compressed). The
variable-length signature is a non-trivial engineering challenge for
block-size accounting (\S\ref{sec:parallel}). Classical security is
$\approx 2^{128}$ operations; there are no known quantum attacks.

Kyber-1024~\cite{fips203} (ML-KEM, FIPS~203) is a lattice-based KEM based
on Module-LWE, providing 256-bit post-quantum security. Parameters: public
key 1568\,B, ciphertext 1568\,B, shared secret 32\,B. In QUANTA, peers
execute a fresh Kyber encapsulation handshake per session to derive a shared
symmetric key. Nodes that do not support Kyber are rejected; no classical
key-agreement fallback exists.

\subsection{Threat Model}

We consider a CRQC adversary targeting: (i)~signature forgery on on-chain
public keys via Shor's algorithm; (ii)~consensus manipulation via forged
block headers; (iii)~cross-protocol and cross-scheme signature replay.
SHA3-256 PoW is considered Grover-resistant with a 128-bit classical
security margin after the square-root speedup.

% ─────────────────────────────────────────────────────────────────────────────
\section{Protocol Hardening for Falcon-512}
\label{sec:hardening}

The gap between what the Falcon reference library guarantees (correct
cryptographic verification) and what consensus requires (deterministic,
malleability-free, replay-resistant verification) is addressed by four
mechanisms. We give an informal security argument for each.

\subsection{Domain-Separated Canonical Signing}

Every transaction signature covers the hash
$\texttt{SHA3-256}(\texttt{``QUANTA\_TX\_V1:''} \| \textit{signing\_bytes})$,
not the raw transaction bytes.

\begin{lstlisting}[caption={Canonical signing hash construction.}]
SIGNING_DOMAIN = "QUANTA_TX_V1:"
canonical_signing_hash(data):
    hasher = SHA3_256.new()
    hasher.update(SIGNING_DOMAIN)
    hasher.update(data)
    return hasher.finalize()   // 32-byte output
\end{lstlisting}

\textbf{Security argument (cross-protocol replay).}
A valid Falcon signature in any other protocol $\Pi'$ covers a hash without
the domain prefix. These hashes are distinct with probability $1-2^{-256}$
by collision resistance of SHA3-256. QUANTA's verifier rejects any signature
that does not verify against the domain-prefixed hash.

\subsection{Pre-Verification Size Guards}

\begin{lstlisting}[caption={Length bounds enforced before cryptographic verification.}]
FALCON512_PUBKEY_BYTES  = 897   // exact; consensus-critical
FALCON512_SIG_MAX_BYTES = 698   // 666-byte sig + 32-byte message
FALCON512_SIG_MIN_BYTES = 33

verify_strict(message, signed_msg, public_key):
    if len(public_key) != FALCON512_PUBKEY_BYTES:  return False
    if not (SIG_MIN <= len(signed_msg) <= SIG_MAX): return False
    // cryptographic verification follows ...
\end{lstlisting}

\textbf{Security argument (DoS hardening).}
Inputs are rejected in $O(1)$ before entering Falcon internals.
The attacker's forcing cost for cryptographic DoS is bounded to
constructing correctly-sized but otherwise invalid inputs, which Falcon's
internal checks reject in bounded time.

\subsection{Sender Binding}

Address derivation computes $\texttt{addr} = \texttt{SHA3-256}(pk)[0..20]$.
The verifier checks this binding \emph{before} calling Falcon verify.

\textbf{Security argument (key-substitution).}
An adversary holding $(sk', pk')$ with $pk' \neq pk_v$ cannot satisfy the
binding check; SHA3-256 collision resistance makes
$\texttt{SHA3-256}(pk')[0..20] = \texttt{SHA3-256}(pk_v)[0..20]$
computationally infeasible.

\subsection{Memory Safety and Scheme Tag}

Secret key storage is zeroed on drop, preventing recovery via heap or swap
inspection. Each transaction carries a 1-byte algorithm tag; unknown tag
values cause the transaction to be hard-rejected. This prevents scheme
downgrade attacks: no valid transaction with an unrecognised or
classical-scheme tag can propagate through a QUANTA node.

% ─────────────────────────────────────────────────────────────────────────────
\section{Parallel and Cached Block Validation}
\label{sec:parallel}

Falcon-512 verification costs approximately 1.5\,ms per operation on a
modern x86-64 core~\cite{pornin2022}. A maximum-capacity block of
1,200~transactions would require $\approx$1,800\,ms serially---well above
the 10-second block budget before balance and nonce checks.

\subsection{Parallel Verification}

Block validation distributes signature verification across all available
CPU cores using a data-parallel work-stealing scheduler.

\begin{lstlisting}[caption={Data-parallel block validation.}]
// Pure-function semantics of Falcon verify make this safe.
all_valid = block.transactions
    .parallel_iter()
    .all(|tx| tx.is_system_tx() OR tx.verify_signature())
\end{lstlisting}

Signature verification is a pure function, so data-parallel evaluation
is safe. Balance and nonce validation, which requires sequential access to
account state, follows in a separate serial pass.

\subsection{LRU Signature Cache and Variable-Length Accounting}

Transactions verified during mempool admission need not be re-verified at
block arrival. A 100,000-entry LRU cache maps each transaction's SHA3-256
hash to its Boolean verification result. Because Falcon-512 signatures are
33--666\,B, block construction serializes each candidate transaction and
accumulates byte counts, skipping any transaction that would push
the block past the 2\,MB limit.

\begin{table}[h]
\centering
\begin{tabular}{@{}lrr@{}}
\toprule
\textbf{Strategy} & \textbf{Projected time (ms)} & \textbf{Basis} \\
\midrule
Serial, cold cache           & $\approx$1{,}800 & $1{,}200 \times 1.5$\,ms \cite{pornin2022} \\
Parallel, 8 threads          & $\approx$225     & ideal $8\times$ speedup   \\
Parallel + 80\% LRU hit      & $\approx$45      & 20\% of 225\,ms           \\
\bottomrule
\end{tabular}
\caption{Projected block-validation latency (1,200\,tx block). Figures are
analytical estimates from published Falcon-512 verification
costs~\cite{pornin2022}, not measured results.}
\label{tab:timing}
\end{table}

Even the conservative serial projection fits within the 10-second block
budget. The parallel and cached configurations provide headroom for the
subsequent sequential passes.

% ─────────────────────────────────────────────────────────────────────────────
\section{Build and Consensus Determinism}
\label{sec:determinism}

If two honest nodes produce different verification results for the same
transaction due to compiler differences, the network forks. QUANTA enforces
determinism at three levels.

\textbf{Toolchain pinning.}
The Rust compiler is pinned to version 1.85.0 via a repository-level
toolchain specification. Nightly or beta channels are prohibited: a different
LLVM back-end may apply different loop-vectorization passes to the Falcon C
library, potentially changing execution timing or floating-point rounding in
auxiliary computations, leading to consensus divergence.

\textbf{Integer-only consensus arithmetic.}
All consensus computations (block reward, difficulty adjustment) use pure
64-bit unsigned integer arithmetic. For example:
\[
  \text{reward} \;\leftarrow\; \lfloor \text{reward} \times 85 \;/\; 100 \rfloor
\]
IEEE~754 double-precision would produce different results on x87 FPU
(80-bit extended precision) vs.\ ARM64 strict-IEEE mode (64-bit), creating
cross-architecture consensus forks. The integer approach accumulates less
than 0.01\,QUA error over 20 years, well within the 5\,QUA minimum reward.

\textbf{Reproducible build profile.}
A single code-generation unit (eliminates non-deterministic LLVM inlining),
link-time optimization (LTO), and disabled incremental compilation ensure
byte-for-byte binary reproducibility across runs on the same platform.

% ─────────────────────────────────────────────────────────────────────────────
\section{Compressed Storage}
\label{sec:storage}

Falcon-512 transactions carry 897-byte public keys and up to 666-byte
signatures in every record. At 1,200~transactions, na\"ive JSON encoding
occupies $\approx$3.2\,MB per block. QUANTA implements a two-stage pipeline:

\textbf{Binary serialization.}
Replacing JSON with a compact binary encoding eliminates field-name overhead,
braces, and ASCII representation of binary data, yielding $\approx$22\%
size reduction and $\approx$8$\times$ throughput improvement.

\textbf{Zstandard compression (level 3).}
After binary serialization, each block is compressed with Zstandard before
being written to the embedded key-value store. Zstandard exploits the high
repetition of identical 897-byte public keys across transactions in the same
block, achieving $\approx$3--4$\times$ compression on blockchain data.
Combined, the two stages reduce block size to $\approx$0.8\,MB---a 75\%
reduction vs.\ plain JSON.

\textbf{On-demand loading and transaction indexing.}
The in-memory blockchain retains only the genesis block; all others are
loaded on demand from the persistent store, reducing RAM usage from $O(h)$
(chain height) to $O(1)$. A 1,000-block LRU cache serves common queries.
A secondary index maps each transaction hash to a
(\textit{block-height}, \textit{position}) tuple, reducing lookup from
$O(h)$ sequential scan to $O(1)$.

% ─────────────────────────────────────────────────────────────────────────────
\section{Crypto-Agility}
\label{sec:agility}

Despite Falcon-512's strong security guarantees, prudent design mandates a
clear upgrade path. QUANTA embeds a single-byte algorithm tag in every
transaction. The dispatcher hard-rejects unknown tag values before any
cryptographic work begins. New algorithms are activated by assigning a new
tag and issuing a soft fork: upgraded nodes accept both old and new tags;
unupgraded nodes reject the new tag but continue to accept the old one.
No wire-format changes are required.

\begin{table}[h]
\centering
\begin{tabular}{@{}cll@{}}
\toprule
\textbf{Tag} & \textbf{Algorithm} & \textbf{Status} \\
\midrule
\texttt{0x01} & Falcon-512           & Active \\
\texttt{0x02} & ML-DSA (Dilithium)   & Reserved \\
\texttt{0x03} & SLH-DSA (SPHINCS+)   & Reserved \\
\bottomrule
\end{tabular}
\caption{Algorithm tag space in the QUANTA wire format.}
\end{table}

The domain separator includes an explicit version token
(\texttt{QUANTA\_TX\_V1:}), ensuring that signatures under distinct scheme
versions cannot be replayed across versions. This design is informed by
NIST's guidance to plan for algorithm migration in long-lived
systems~\cite{fips206}.

% ─────────────────────────────────────────────────────────────────────────────
\section{Implementation}

QUANTA is implemented in $\approx$9,000 lines of Rust (2021 edition).
Rust's ownership model eliminates memory-safety vulnerabilities at compile
time and statically prevents data races across the parallel verification
threads.

\begin{table}[h]
\centering
\begin{tabular}{@{}ll@{}}
\toprule
\textbf{Component} & \textbf{Library / Standard} \\
\midrule
Async runtime        & Tokio 1.35 \\
PQ signatures        & pqcrypto-falcon 0.3.0 (version-pinned) \\
PQ key encapsulation & pqcrypto-kyber 0.8 (Kyber-1024) \\
Hashing              & SHA3-256 (FIPS 202) \\
Storage              & sled 0.34 (embedded key-value store) \\
Compression          & Zstandard 0.13 \\
Parallel compute     & Work-stealing thread pool \\
Verification cache   & LRU cache, 100k entries \\
\bottomrule
\end{tabular}
\caption{QUANTA runtime stack.}
\end{table}

The Falcon library dependency is pinned to an exact version; any version
change alters the compiled C code and could change signature encoding,
requiring a coordinated network-wide upgrade. HD wallet derivation uses
BIP39 24-word mnemonics with Argon2id key stretching; wallet files encrypt
Falcon-512 secret keys with ChaCha20-Poly1305.

The code is MIT-licensed at \url{https://github.com/quantachain/quanta};
a pre-mainnet testnet node is available via Docker
(\texttt{docker pull xd637/quanta-node:latest}).
A systematic throughput benchmark study over a multi-operator testnet is
in progress and will be reported in a revision of the extended
preprint~\cite{quanta2026}.

% ─────────────────────────────────────────────────────────────────────────────
\section{Discussion and Related Work}

\textbf{QRL}~\cite{hulsing2018} uses XMSS (RFC~8391), whose stateful
signing requires per-key state. This significantly complicates wallet
management and key recovery in a blockchain context.

\textbf{QAN Platform} deploys Kyber and Dilithium on an EVM-compatible
layer-1. EVM account model complexity is inherited, and cross-architecture
consensus determinism is not addressed.

\textbf{Algorand's PQ research} pursues hybrid ECDSA+Falcon signatures.
Retaining ECDSA in the critical path means the hybrid construction offers
only the security of the weaker primitive against a CRQC.

\textbf{Ethereum's roadmap}~\cite{buterin2023} acknowledges eventual PQ
migration without committing to a concrete algorithm or timeline, leaving
all existing account keys at risk once CRQCs arrive.

\textbf{Limitations.}
Proof-of-Work carries well-known environmental concerns; future versions
may explore proof-of-stake with post-quantum verifiable random functions.
Falcon-512 signatures (up to 666\,B) are larger than ECDSA (71\,B) or
Ed25519 (64\,B); the compression layer partially offsets this, and the
crypto-agility design makes migration to a more compact scheme feasible.
The current deployment is a pre-mainnet testnet; production readiness
requires external security audits and extended stress testing.

% ─────────────────────────────────────────────────────────────────────────────
\section{Conclusion}

We presented QUANTA, an open-source blockchain built exclusively on
NIST-standardized post-quantum cryptography. The five engineering
contributions address the gap between academic PQ scheme correctness and
the requirements of a production consensus system: domain-separated
verification prevents cross-protocol replay; parallel verification and LRU
caching make 1,200-transaction block validation feasible within a 10-second
block budget; integer-only arithmetic eliminates floating-point
non-determinism across architectures; Zstandard compression reduces on-disk
block size by 75\%; and the single-byte agility tag enables fork-free
algorithm migration. QUANTA demonstrates that post-quantum blockchains are
practically engineerable today on commodity hardware.

% ─────────────────────────────────────────────────────────────────────────────
\begin{thebibliography}{10}

\bibitem{shor1994}
P.~W. Shor,
``Algorithms for quantum computation: discrete logarithms and factoring,''
\textit{Proc.\ 35th FOCS}, pp.~124--134, 1994.

\bibitem{mosca2018}
M.~Mosca,
``Cybersecurity in an era with quantum computers: will we be ready?''
\textit{IEEE Security \& Privacy}, vol.~16, no.~5, pp.~38--41, 2018.

\bibitem{hulsing2018}
A.~H\"ulsing et al.,
``XMSS: Extended Merkle Signature Scheme,'' RFC~8391, IETF, 2018.

\bibitem{falcon2020}
P.-A. Fouque et al.,
``Falcon: Fast-Fourier Lattice-based Compact Signatures over NTRU,''
NIST PQC Round~3 Submission, 2020. \url{https://falcon-sign.info/}

\bibitem{fips203}
NIST,
``Module-Lattice-Based Key-Encapsulation Mechanism Standard,''
FIPS~203, 2024.

\bibitem{fips206}
NIST,
``Module-Lattice-Based Digital Signature Standard,''
FIPS~206, 2024.

\bibitem{buterin2023}
V.~Buterin,
``The Ethereum roadmap: post-quantum migration,''
Ethereum Research Forum, 2023. \url{https://ethresear.ch/}

\bibitem{pornin2022}
T.~Pornin,
``New Efficient, Constant-Time Implementations of Falcon,''
\textit{IACR TCHES}, 2022.

\bibitem{quanta2026}
K.~K.,
``QUANTA: Engineering a Production-Ready Post-Quantum Blockchain
with Falcon-512 Lattice Signatures,''
Zenodo, Feb.\ 2026.
DOI: \href{https://doi.org/10.5281/zenodo.18753528}{10.5281/zenodo.18753528}.

\end{thebibliography}

\end{document}
